\uuid{KoeL}
\exo7id{1639}
\titre{exo7 1639}
\auteur{barraud}
\organisation{exo7}
\datecreate{2003-09-01}
\isIndication{false}
\isCorrection{false}
\chapitre{Réduction d'endomorphisme, polynôme annulateur}
\sousChapitre{Diagonalisation}
\cpge{mpsi-mp2i.calcul-matriciel-systemes.e, pcsi.determinants.c, ptsi.determinants.a, mpsi-mp2i.changements-bases.a, pcsi.matrices-applications-lineaires.c, ptsi.matrices-applications-lineaires.c, bcpst1.matrices, mp-mpi.reduction-endomorphismes.d, pc.reduction.c, psi.reduction-endomorphismes.c, pt.reduction-endomorphismes.c, mp-mpi.endomorphismes-euclidien.b, pc.endomorphismes-euclidiens.b, psi.endomorphismes-euclidien.b, pt.isometries-euclidien.b, mp-mpi.endomorphismes-euclidien.f, pc.endomorphismes-euclidiens.d, psi.endomorphismes-euclidien.f, pt.isometries-euclidien.e, bcpst2.valeurs-propres-vecteurs-propres.b, bcpst2.produit-scalaire.c}
\module{Algèbre}
\niveau{L2}
\difficulte{}

\contenu{
\question{
Soit $A$ la matrice $A=
  \begin{pmatrix}
    1&-1&-1\\
   -1& 1&-1\\
   -1&-1& 1
  \end{pmatrix}
  $. Calculer $^{t}\!A$. La matrice $A$ est-elle diagonalisable~?

  Trouver une matrice $P$ orthogonale telle que $P^{-1}AP$ soit
  diagonale.
}
}
